\documentclass[11pt]{article}
% fontenc genera PDF con caracteres latinos copiables.
\usepackage[T1]{fontenc}
% inputenc permite que el archivo fuente esté codificado en UTF-8.
\usepackage[utf8]{inputenc}
% babel adapta al español los nombres y la separación silábica.
\usepackage[spanish]{babel}
% amsmath proporciona los entornos matemáticos habituales.
\usepackage{amsmath}
% amssymb aporta los conjuntos numéricos \mathbb del tablero de símbolos.
\usepackage{amssymb}
% amsthm permite declarar teoremas y bloques relacionados.
\usepackage{amsthm}
% Los entornos siguientes se derivan de la tabla de tipos de bloque.
\theoremstyle{plain}
\newtheorem{theorem}{Teorema}
\newtheorem{proposition}{Proposición}
\theoremstyle{definition}
\newtheorem{definition}{Definición}
\newtheorem{example}{Ejemplo}
\theoremstyle{remark}
\newtheorem{note}{Nota}

\title{Notas de Cálculo III}
\author{Ana Pérez \\ Cálculo III \\ Profesor: Dr. Ruiz}
\date{2026-09-11}

\begin{document}

\maketitle

\section{Integrales múltiples}

\begin{definition}[Integral doble]
Sea $f: A \to \mathbb{R}$ acotada en $A \subseteq \mathbb{R}^2$. La integral sobre $A$ se escribe en la ecuación siguiente.
\end{definition}

\[
\iint_A f(x,y) \, dx \, dy
\]

\begin{theorem}[Fubini]
Bajo las hipótesis del curso, el orden de integración no altera el resultado: se cumple $\iint_A f = \int \! \int f \, dx \, dy$ siempre que
\begin{itemize}
\item $f$ sea continua en el rectángulo $[a,b] \times [c,d]$
\item el 100\% del recinto quede dentro de $A$ \& sin cortes
\end{itemize}
\end{theorem}

\begin{example}[Rectángulo]
Para $f(x,y)=x+y$ en $[0,1] \times [0,2]$, calculamos el valor por iteración con \textbf{el orden natural}.

Este bloque tiene dos párrafos y llega al .tex tal cual se escribió.
\end{example}

\begin{note}
La argumentación escrita cuenta para la calificación: no basta con el símbolo.
\end{note}

\end{document}
